\documentclass[10pt,a4paper]{amsart}
\usepackage[margin=19mm]{geometry}
\usepackage{amsmath,amssymb,mathtools}
\usepackage[scr=rsfs,cal=euler]{mathalfa}
\usepackage{xfrac}
\usepackage[unicode,hidelinks,bookmarks=false]{hyperref}
\newcommand{\ie}{i.e.\ }
\newcommand{\cf}{cf.\ }
\newcommand{\resp}{respectively\ }
\newcommand{\Subsection}{\subsection}
\providecommand{\nobreakdash}{\nobreak\hbox{-}}
\setlength{\emergencystretch}{2em}
\multlinegap0pt
\usepackage{bm}
\usepackage{bbm}
\usepackage{dsfont}
\usepackage{xspace}
\usepackage{cancel}
\usepackage{mathtools}
\usepackage{amsthm}
\makeatletter
\@ifundefined{theorem}{\newtheorem{theorem}{Theorem}}{}
\@ifundefined{lemma}{\newtheorem{lemma}[theorem]{\bf Lemma}}{}
\makeatother
\title[Extending edge colorings of distance-3 matchings in the Cart]{Extending edge colorings of distance-3 matchings in the Cartesian product of graphs}
\author{P\'al B\"arnkopf, Ervin Gy\H{o}ri}
\date{Source: arXiv:2310.09973v6 (CC BY 4.0)}
\begin{document}
\maketitle

\noindent\emph{Source and licence.} Extending edge colorings of distance-3 matchings in the Cartesian product of graphs. Version arXiv:2310.09973v6 (CC BY 4.0), distributed under the Creative Commons Attribution 4.0 International licence, as recorded in the arXiv licence metadata for this exact version and shown on the abstract page https://arxiv.org/abs/2310.09973v6. Full rights notice: licenses/arxiv-2310.09973-CC-BY-4.0.txt.\par\smallskip

\noindent\textit{Editorial scope.} Counted unit: the Lemma whose source label is \texttt{regularization} in the article (source0.tex lines 227--229, source label \texttt{regularization}), delivered together with the article's own complete original proof of it (source0.tex lines 231--249, one \texttt{proof} environment of 19 lines). No mathematics is written, repaired or re-derived: statement and proof are the article's own text line for line. Required context retained uncounted: none. Every definition, display and label of those lines is retained unchanged. Omitted from this package and not redistributed: 1--226, 230--230, 250--540. Nothing inside a retained excerpt is omitted or altered: every retained body is delivered exactly as the article's lines are written, so no omission receipt of any kind accompanies this package. Citations: the retained excerpts carry no citation command at all, so no bibliography entry of the article is transcribed and no reference is redistributed. Reference adaptation: none was needed and none was made; every reference inside the delivered unit resolves inside it, so no unresolved reference or citation is produced. Printed slips kept literal: no printed word, symbol, hypothesis or number of the counted statement or of its proof was altered. Layout and class: the article's own class is \texttt{article} with packages including amsmath, amsthm, mathtools, amssymb, amscd, epic, eepic, url, color, inputenc, comment, todonotes, graphicx, epstopdf; the wrapper is the lane's pinned \texttt{amsart} editorial wrapper at 10pt on A4, which restates the theorem-like environments the retained text needs and adds the article's own macro definitions that the retained text needs (none), with extra wrapper packages (bm, bbm, dsfont, xspace, cancel, mathtools). The delivered numbering is the unit's own. Assets: no retained span contains a figure environment, a graphics-inclusion command, a PDF-page inclusion, a table or a TikZ picture; no image file is copied into this package, the package declares no asset dependency and no image file is redistributed with it. Licence: version arXiv:2310.09973v6 of this article is distributed under the Creative Commons Attribution 4.0 International licence, as recorded in the arXiv licence metadata for this exact version and shown on the abstract page https://arxiv.org/abs/2310.09973v6. A full rights notice is delivered at licenses/arxiv-2310.09973-CC-BY-4.0.txt. Characters: every retained excerpt is pure ASCII, so the delivered engine, XeLaTeX with its default Latin Modern text font, reports no missing character.
\begin{lemma} \label{regularization}
    Let $k$ be an integer with $k \geq 2$, and let each $G_i$ (for $i \in \{1, \dots, k\}$) be a Class 1 graph. By adding extra vertices and edges to each graph $G_i$, we can make them $\Delta(G_i)$-regular while preserving their Class 1 property and ensuring that if the distance between any two edges in the Cartesian product of the original graphs $G_i$ is at least 3, then the same holds in the Cartesian product of the regularized version as well.
\end{lemma}

\begin{proof}
    It suffices to show that a single factor can be made regular. This procedure can then be applied individually to each factor, thereby proving the statement.

    We describe a possible procedure for adding these extra vertices and edges to a graph $H$. Starting with $H$, whenever a vertex has a degree less than $\Delta(H)$, we duplicate the entire graph and connect each such vertex to its counterpart in the copy. This ensures that any vertex of degree less than $\Delta(H)$ gains exactly one new neighbor, increasing its degree without exceeding $\Delta(H)$. This operation increases the minimum degree without raising the chromatic index. We repeat this process until the graph becomes regular. (The minimum degree increases in every step and the maximum degree is unchanged, so the process stops in $\Delta(H)-\delta(H)$ steps.)

    We claim that the graph remains Class 1 at every step. $H$ has a $\Delta(H)$-edge-coloring and we can color both copies identically. Furthermore, new edges are only added between vertices whose degrees are still below $\Delta(H)$, so at each such vertex there is an unused color. We can assign these unused colors to these new edges without increasing the total number of colors used.

    Now, let $G$ and $H$ be Class 1 graphs and let $H^*$ denote the regular extension of $H$ obtained by this process. We show that if the distance between two edges $e_1, e_2$ in $G \square H$ is at least 3, it remains at least 3 in $G \square H^*$. Assume, for contradiction, that $d_{G \square H}(e_1, e_2) \geq 3$ but $d_{G \square H^*}(e_1, e_2) \leq 2$.

    Since no new edges were added between original vertices, the distance cannot be 1. (And the distance cannot be 0, of course.) If the distance is 2, there must exist a vertex $(a, b)$ in $G \square H^*$ and vertices $(u, v) \in e_1$ and $(x, y) \in e_2$ such that $(u, v) \sim (a, b)$ and $(a, b) \sim (x, y)$. (We  use $\sim$ to indicate adjacency.)
    
    \begin{enumerate}
        \item \textbf{Case 1: $(a, b)$ is an original vertex in $G \square H$.} Since the construction only adds edges between original and new vertices (or between two new vertices), the adjacency relations between original vertices are unchanged. This implies $d_{G \square H}((u, v), (x, y)) = 2$, contradicting $d_{G \square H}(e_1, e_2) \geq 3$.
        
        \item \textbf{Case 2: $(a, b)$ is a new vertex.} By the definition of the Cartesian product, for the adjacency $(u, v) \sim (a, b)$ where $(u, v)$ is original and $(a, b)$ is new, we must have $u = a$ and $v \sim  b$ in $H^*$. Similarly, for $(x, y) \sim (a, b)$, we must have $x = a$ and $y \sim  b$ in $H^*$. In our regularization procedure, whenever we duplicate a graph, we only add edges between a vertex and its direct copy (or between two new vertices). By induction, any new vertex $b \in V(H^*) \setminus V(H)$ is connected to at most one vertex in the original vertex set $V(H)$, which is its unique "preimage" (the vertex from which it was originally derived). Since $v \sim  b$ and $y \sim  b$ in $H^*$, and $v, y \in V(H)$ while $b \notin V(H)$, the construction implies that $v$ and $y$ both must be the unique preimage of $b$ in $V(H)$. It follows that $u = x$ and $v = y$. Thus, $(u, v) = (x, y)$, which means that $e_1$ and $e_2$ share a vertex, so their distance is 0, a contradiction.
    \end{enumerate}    

    Since all cases lead to a contradiction, our assumption was false: the distance cannot decrease below 3 in the regularized version.
\end{proof}

\end{document}
